% Section 1 overview
A smooth solution history begins locally, but a finite maximal time remains possible unless the equations keep one fixed restart-grade Sobolev norm finite.
The known estimates narrow that failure without closing it, leaving a direct Gold obligation and its logically equivalent Silver obligation.
Both obligations concern the same transported velocity--pressure history, so their meaning has to be tested by asking whether a proposed continuation still belongs to the original problem, still participates in the surrounding field, and still carries every required response of that one fluid.

% Silver ambient universe
For this explanatory picture, let \(\mathcal Q\) be the collection of proposed continuations for which the preceding same-history, distributional, energy, decay, and pressure-normalization requirements are meaningful.
This is a schematic universe for the Silver comparison; the closure in \cref{sec:FromCriticalHeightToCompatibleFiniteViews,sec:TheCompatibleIntersectionAtTheTerminalTime} follows the fixed pressure-normalized maximal classical history generated by \((u_0,\nu)\), not an arbitrary \(Q\in\mathcal Q\).
Within this picture, define

% Section 1.9 opening
Restart sharpens the one-fluid burden.
A terminal slice can restart the same history only if velocity, pressure, and every required response arrive there as one compatible field.
Call that joined relation the VPI law: viscosity, pressure, and incompressibility acting through one transported velocity.
Dead, Jump, and Blown name ways that a portion of the field can lose this relation with the fluid around it.
Following a material packet makes the relation literal.

% Section 1.9 formal VPI lead
The preceding relations are simultaneous parts of the VPI law.
In compact form, viscosity, pressure, and incompressibility act through one transported velocity:

% Section 1.9 moving-neighbour replacement
Take one parcel together with the nearby parcels carried beside it by the exact flow map.
As their distances and directions change, viscosity reads the curvature of their shared velocity neighbourhood, pressure responds to the whole divergence-free field containing them, and the resulting acceleration remains the next response of that same parcel.

% Section 2 overview
The one-fluid account leaves a finite-endpoint threat that can hide in smaller scales while kinetic energy remains finite.
Critical height keeps that concentration visible, but finite time forces an escaping response to call on the response immediately above it, returning temporal depth to the spatial and pressure structure of the original equation.
Finite adjacent rectangles join neighbouring temporal and spatial rows at common traces.
For these rectangles to reach one terminal state, the family must be compatible and each fixed rectangle must have a cutoff-uniform terminal bound.

% Section 1.8 duplicate Zeno removal

% Section 2.1 Zeno callback
The finite parabolic schedule already developed in \cref{sub:HistoricalEstimates} shows that progressively smaller scale durations can fit before one finite time.
At the critical height just derived, that schedule remains dimensionally possible while the \(L^2\) contribution shrinks, but it does not show that the equation generates the successive states.
The actual fluid would still have to create each finer state from the state before it while transport, pressure, incompressibility, and viscosity all continue to act.

% Section 2.2 material callback
The material ladder was built in \cref{sub:SingularityAndTheDerivativeTower}: along one flow line, \(A_r=D_t^ru(X,t)\) and \(dA_r/dt=A_{r+1}\).
Its fundamental-theorem identity gives the conditional mechanism needed here: if \(A_r\) escapes by \(T_*\), then \(A_{r+1}\) cannot remain integrable along that path.
This does not make any rung monotone or nonzero; it states what finite-time escape would require inside the same fluid.

% Section 2.2 handoff after fixed-B escalation
The material and Eulerian ladders are not identical, as \cref{sub:SingularityAndTheDerivativeTower} already showed: each material derivative brings transport into the temporal rows and mixes them with spatial derivatives and nonlinear products.
The particle picture supplies the escalation, but the equation must keep the temporal, spatial, and pressure pieces coupled.
A single time index can no longer hold the response, so the next object must retain the full mixed jet.

% Section 2.3 callback
The pressure-complete temporal recurrence in \cref{sub:SingularityAndTheDerivativeTower} already couples \(V_{n+1}\) to \(\nu\Delta V_n\), the binomial transport sum, and \(\nabla Q_n\), while the normalized pressure Poisson row recovers \(Q_n\) from those same velocity rows.
The fixed-space escalation above explains why adjacent temporal rows matter near a finite endpoint.
What the temporal notation still hides is how spatial derivatives distribute through each nonlinear product.
